\documentclass[letterpaper,10pt,conference]{ieeeconf}
\IEEEoverridecommandlockouts
\overrideIEEEmargins

\usepackage{amsmath,amssymb,bm}
\usepackage{booktabs}
\usepackage{graphicx}
\usepackage{multirow}
\usepackage{url}
\usepackage{cite}
\usepackage{xcolor}
\usepackage{tikz}
\usetikzlibrary{arrows.meta,positioning,fit}

% Citation/link blue sampled from the CoSAM manuscript (#367DBD).
\definecolor{cosamLinkBlue}{HTML}{367DBD}

% ieeeconf defines an obsolete Natbib compatibility hook.  Removing it before
% loading hyperref restores clickable numerical citations.
\makeatletter
\let\NAT@parse\undefined
\makeatother
\usepackage[
  colorlinks=true,
  linkcolor=cosamLinkBlue,
  citecolor=cosamLinkBlue,
  urlcolor=black
]{hyperref}

\definecolor{todoRed}{RGB}{176,36,36}
\definecolor{inputBlue}{RGB}{225,238,250}
\definecolor{geometryGreen}{RGB}{224,243,231}
\definecolor{routeOrange}{RGB}{253,237,214}
\definecolor{outputPurple}{RGB}{238,229,247}
\newcommand{\todo}[1]{\textcolor{todoRed}{\textbf{[TODO: #1]}}}
\newcommand{\tbd}{\textcolor{todoRed}{--}}
\newcommand{\R}{\mathbb{R}}
\newcommand{\I}{\mathbf{I}}
\newcommand{\D}{\mathbf{D}}
\newcommand{\K}{\mathbf{K}}
\newcommand{\T}{\mathbf{T}}
\newcommand{\M}{\mathbf{M}}
\newcommand{\F}{\mathbf{F}}
\newcommand{\U}{\mathbf{U}}
\newcommand{\n}{\mathbf{n}}
\newcommand{\od}{\mathbin{\odot}}
\newsavebox{\roboticplatformbox}

\title{\LARGE \bf
DistriNVS: Distributional Geometry Transport and Visibility-Constrained Completion for Cross-Endoscope Novel-View Synthesis}

\author{Anonymous Authors}

\begin{document}
\maketitle
\thispagestyle{empty}
\pagestyle{empty}

\begin{abstract}
Cross-endoscope novel-view synthesis must both transport tissue observed by
the source camera and infer tissue revealed only from the target viewpoint.
These two regimes are confounded when sparse or erroneous depth is projected
as deterministic points: sampling gaps, depth failures, surface collisions,
and genuine disocclusions all become visually identical holes. We present
DistriNVS, a reliability-aware framework that models every source depth as a
distribution and propagates it through calibrated geometry. Its
distributional surface splatting renderer produces an anisotropic target-plane
footprint together with support, depth variance, and collision entropy. These
statistics drive reliability-constrained transport--synthesis fusion: a bounded
residual expert refines transported appearance, whereas an
uncertainty-conditioned Fourier network predicts unsupported or geometrically
unreliable content. Their reliability-weighted RGB-D mixture and an explicit
risk head are trained with target-depth and backward-reprojection constraints.
The prescribed protocol trains on a fixed
set of 13 iMED scenes and evaluates without adaptation on held-out EndoVis
datasets 8 and 9 using a strict raw-overlap protocol. 
% \todo{Insert the final test-set
% sample count, the strongest quantitative improvements with confidence
% intervals, and runtime after all experiments are frozen.}
\end{abstract}

\section{Introduction}
\label{sec:introduction}

Robot-assisted minimally invasive surgery depends on endoscopic vision for
scene perception and instrument guidance
~\cite{wei2026roboticperception,sunmola2026surgicalgaussiansurfels}.
Cross-endoscope novel-view synthesis can broaden this visual context by
predicting a target-endoscope view from synchronized appearance and geometry
acquired by a spatially displaced source endoscope~\cite{wang2026endoe2egs}.
Unlike left-to-right synthesis within one stereo endoscope, the target may
observe anatomy absent from both source cameras. The task therefore combines
two epistemically distinct regimes: co-visible anatomy has measured appearance
that should be transported through calibrated geometry, whereas target-only
anatomy lacks source evidence and must be synthesized. A reliable system should
retain source evidence, identify inferred content, and thereby support
wider-field robotic perception, retrospective visualization, and surgical
simulation~\cite{hsieh2026extrags}.
Such uncertainty-aware cross-view completion may also support visual-state
construction for surgical robot world models~\cite{he2025surgworld}.

Achieving this separation is difficult because endoscopic depth and calibration
are unreliable around specularities, weak texture, depth discontinuities,
instruments, and deforming tissue~\cite{allan2021scared}. Conventional forward
warping represents each valid source pixel by one 3-D point; discretized
projection leaves sampling gaps, while multiple points may compete for the same
target pixel. A binary validity mask then conflates missing depth, rasterization
holes, depth-ordering conflicts, and genuine cross-endoscope disocclusions.
Treating these cases uniformly either degrades observed structure or permits
hallucination where geometric evidence remains trustworthy.

\begin{figure}[t]
\centering
\vspace*{2mm}
\sbox{\roboticplatformbox}{%
  \includegraphics[width=\columnwidth]{figures/distrinvs_teaser/distrinvs_teaser_editable.pdf}%
}
\begin{tikzpicture}
  \clip[rounded corners=0mm]
    (0,0) rectangle
    (\wd\roboticplatformbox,\ht\roboticplatformbox);
  \node[anchor=south west,inner sep=0pt] at (0,0)
    {\usebox{\roboticplatformbox}};
\end{tikzpicture}
\caption{Cross-endoscope novel-view synthesis under partial field-of-view
overlap. Given the calibrated source-to-target pose, co-visible anatomy is
source-supported and can be transported geometrically (\textcolor{green}{green}), whereas
target-only anatomy lacks source evidence and must be synthesized (\textcolor{purple}{purple}).
DistriNVS represents this provenance boundary through reliability-conditioned
routing between transported and synthesized RGB-D.}
\label{fig:conditional_refinement}
\vspace{-0.4cm}
\end{figure}

Existing methods address complementary parts of this problem but do not expose
the required evidence boundary. Differentiable point rendering and softmax
splatting improve feature transport and collision handling
~\cite{wiles2020synsin,niklaus2020softmax}, while reference-guided inpainting
combines depth-based alignment with learned completion~\cite{zhao2023geofill}.
Large-mask inpainting networks provide strong global image priors
~\cite{suvorov2022lama,chu2023uffc}, and endoscopic video completion adds
depth, temporal, or stereo cues
~\cite{daher2023temporal,zhang2024daevi,zou2025ssifnet}. Endoscopic neural
rendering has likewise advanced sparse-view reconstruction and uncertainty
modeling~\cite{li2024endosparse,guo2025ucnerf}. However, deterministic
projection commonly assigns equal reliability to all valid depths, while
image-space uncertainty does not explicitly describe how source geometry
supports a spatially offset target view. Neither approach propagates per-pixel
depth uncertainty through the cross-endoscope transform and uses the resulting
visibility statistics to constrain where appearance may be synthesized.

Our central premise is to \emph{transport reliable geometric evidence and
synthesize only unsupported or geometrically unreliable regions}. We introduce
DistriNVS, a reliability-aware framework
for direct cross-endoscope novel-view synthesis. Its distributional surface
splatting (DSS) represents each source depth as a distribution whose bounded
mean and scale, surface normal, and confidence are predicted by a depth
reliability encoder. DSS projects deterministic sigma samples through the
calibrated cross-endoscope transform and aggregates source features using
Jacobian- and normal-aware footprints with soft nearest-surface visibility,
yielding target support, depth variance, and collision entropy. These statistics
constrain a soft mixture between bounded residual transport for co-visible
regions and uncertainty-conditioned Fourier completion for unsupported or
ambiguous regions. Target-depth and risk prediction plus backward RGB-D
reprojection constrain this mixture during training.

The paper makes three contributions:
\begin{itemize}
    \item We formulate cross-endoscope novel-view synthesis as a
    reliability-aware separation between transport of anatomy observed by the
    source endoscope and constrained inference of content revealed only by the
    target endoscope.
    \item We introduce DSS, a differentiable renderer that propagates depth
    distributions through calibrated cross-endoscope geometry and exposes
    target support, depth variance, and collision entropy beyond binary validity.
    \item We design a two-expert decoder that softly mixes bounded residual
    transport with uncertainty-conditioned Fourier completion under a
    renderer-derived reliability constraint.
\end{itemize}

\section{Related Work}
\label{sec:related}

We briefly review image completion, geometry-guided view transport, and
endoscopic neural rendering, with particular attention to how existing
approaches represent missing content, exploit calibrated geometry, and quantify
uncertainty.

\begin{figure*}[t]
\centering
\includegraphics[width=\textwidth]{figures/distrinvs_framework/distrinvs_framework.pdf}
\caption{DistriNVS framework. A reliability encoder predicts source geometry
and features. DSS transports deterministic depth hypotheses through the
calibrated cross-endoscope transform and jointly produces warped evidence and
target-plane visibility statistics. These statistics condition a
reliability-constrained soft mixture of a residual transport expert and a
Fourier synthesis expert. Solid arrows denote inference; the dashed lower lane is used only
during training.}
\label{fig:overview}
\end{figure*}

\subsection{Image and Endoscopic Video Inpainting}

Generic image inpainting combines mask-aware feature extraction with
increasingly broad appearance priors. Partial and gated convolutions prevent
invalid activations from contaminating local features
~\cite{liu2018partial,yu2019gated}, whereas LaMa and UFFC use Fourier operations
for global context and MAT restricts attention to valid tokens
~\cite{suvorov2022lama,chu2023uffc,li2022mat}. Recent diffusion models such as
BrushNet and TurboFill incorporate masked-image features into pretrained
generators or adapt few-step sampling for efficient completion
~\cite{ju2024brushnet,xie2025turbofill}. Despite strong appearance priors, these
methods define holes in a fixed image coordinate system and do not describe
geometric provenance. Endoscopic completion adds domain cues: temporal learning
removes specularities, DAEVI combines spatio-temporal RGB and depth, and SSIFNet
fuses flow and stereo-temporal information for instrument occlusions
~\cite{daher2023temporal,zhang2024daevi,zou2025ssifnet}. Such methods exploit
neighboring frames or views but still restore defects defined in the original
camera coordinates. Cross-endoscope synthesis instead creates holes through
viewpoint transfer: target-only anatomy is unobserved, while rasterization gaps
or depth conflicts may retain source support. DistriNVS therefore uses
calibrated projection statistics to distinguish true disocclusions from
rasterization gaps and depth conflicts.

\subsection{Geometry-Guided Transport and Completion}

Geometry-guided view synthesis preserves correspondence by transporting source
evidence before generating missing content. SynSin predicts and projects a
feature point cloud, softmax splatting resolves many-to-one mappings with
learned importance, and GeoFill refines depth and pose before reprojection and
inpainting~\cite{wiles2020synsin,niklaus2020softmax,zhao2023geofill}. Recent
feed-forward Gaussian renderers use stronger explicit representations: MVSplat
localizes Gaussians with plane-sweep cost volumes, whereas DepthSplat combines
multi-view correspondence with monocular depth priors
~\cite{chen2024mvsplat,xu2025depthsplat}. They reduce per-scene optimization but
infer a radiance representation from multiple calibrated observations, without
an explicit per-pixel boundary between measured and synthesized appearance.
Generative methods target wider gaps. LeftRefill frames reference-conditioned
synthesis as paired-canvas completion, MultiDiff conditions video diffusion on
depth-warped reference views, and point-guided latent diffusion uses
stereo-derived clouds for laparoscopic view synthesis
~\cite{cao2024leftrefill,muller2024multidiff,tang2025pointguided}. These models can
generate plausible structures in large disocclusions, yet plausibility alone
does not reveal which pixels remain geometrically supported. DistriNVS instead
transports a calibrated source observation with distributional splatting,
derives support, variance, and collision statistics during rasterization, and
uses them to constrain soft fusion between transported and synthesized content.

\subsection{Endoscopic Novel-View Synthesis and Uncertainty}

Endoscopic neural rendering must handle non-rigid tissue, view-dependent
illumination, tool occlusions, and sparse camera coverage. EndoNeRF represents
deformation with a neural radiance field, while EndoSurf combines dynamic
modeling with a neural surface~\cite{wang2022endonerf,zha2023endosurf}.
EndoSparse introduces geometric and appearance priors into Gaussian splatting,
and Endo-GSMT uses depth, dense-displacement, and motion priors for dynamic
reconstruction~\cite{li2024endosparse,gou2025endogsmt}. GastroNVS further shows
that large viewpoint changes in real gastroscopy amplify geometric artifacts
and inter-view illumination mismatch~\cite{minai2026gastronvs}. These systems
build persistent scene representations from multiple frames and principally
address reconstruction fidelity within a captured scene. Uncertainty-aware
approaches address sparse evidence from another direction: UC-NeRF conditions a
radiance field on multi-view uncertainty, whereas ExtraGS uses uncertainty to
sample virtual cameras and weight diffusion-refined pseudo-observations for
extrapolation~\cite{guo2025ucnerf,hsieh2026extrags}. Thus, neither uncertainty nor
generative priors are new to endoscopic rendering. The distinction is
operational: DSS propagates a distribution for each input depth through the
calibrated cross-endoscope transform and derives target support, depth variance,
and collision entropy from the same rasterization that transports appearance.
DistriNVS uses these quantities to route completion during unseen-sequence
inference without target-scene optimization.

\section{Method}
\label{sec:method}

\subsection{Problem Formulation and Overview}

After resizing, both endoscope views have spatial size $H\times W$; we omit the
batch dimension throughout. DistriNVS receives source RGB
$\I_s\in[0,1]^{3\times H\times W}$, sparse metric depth
$\D_s\in\R_+^{1\times H\times W}$, and validity mask
$\M_s\in\{0,1\}^{1\times H\times W}$. Calibration comprises intrinsics
$\K_s,\K_t\in\R^{3\times3}$ and a rigid transform
$\T_{s\rightarrow t}\in\mathrm{SE}(3)\subset\R^{4\times4}$. The model predicts
target RGB $\widehat\I_t\in[0,1]^{3\times H\times W}$, depth
$\widehat\D_t\in\R_+^{1\times H\times W}$, and risk
$\U_t\in[0,1]^{1\times H\times W}$. Target RGB-D and the target instrument mask
are used only for training or evaluation.

Figure~\ref{fig:overview} summarizes the inference pipeline. A
reliability encoder estimates source geometry and uncertainty. Distributional
surface splatting (DSS) then transports source evidence and measures target-view
support and ambiguity. These measurements route complementary transport and
synthesis experts, whose RGB-D predictions are combined by a
reliability-constrained soft gate. The backward photometric cycle and risk objective
are used only during training.

\subsection{Reliability-Aware Source Geometry}

Naive depth completion erases the distinction between measured and inferred
geometry. The reliability encoder instead processes source RGB,
median-normalized depth, and its validity mask. It produces a feature map
$\F_s\in\R^{d_f\times H\times W}$ and scalar fields
$\mu,\sigma,c\in\R^{1\times H\times W}$. They represent depth mean and scale
in millimetres, and normalized confidence. It also predicts a unit orientation field
$\n\in\R^{3\times H\times W}$. The normalization scale $s_D$ is a positive
scalar in millimetres, chosen as the median valid source depth, and missing
inputs are set to zero. During training, $\M_s$ is the post-dropout input mask,
whereas $\M_s^{\rm gt}$ has the same dimensions and retains the pre-dropout
supervision.

Measured depths receive a bounded residual correction; missing depths use a
positive completion branch. Both obey the physical range, but disjoint
confidence intervals keep inferred geometry less authoritative than
measurements. The optical-axis-stabilized orientation proxy is learned only
through downstream objectives. It serves as a footprint cue rather than a
metrically supervised normal. Fixed ranges and clipping operations appear in
Sec.~\ref{sec:implementation}.

\subsection{Distributional Surface Splatting}

DSS concatenates source RGB and features as $\mathbf A_s=[\I_s,\F_s]$ and
transports it to a target tensor $\mathbf A_t^w$. Both tensors contain $3+d_f$
channels over the $H\times W$ grid. DSS simultaneously predicts expected depth
$\bar D_t\in\R_+^{1\times H\times W}$. Support $S$, coverage
$C$, depth variance $V$, collision entropy $\mathcal H$, confidence $\kappa$,
and render validity $M^{\rm dss}$ are $1\times H\times W$ target-grid maps.
Support and variance are non-negative; coverage, collision entropy, and
confidence lie in $[0,1]$; validity is binary. All arise from one rasterization
of a compact depth distribution at each source pixel. Expected depth is in
millimetres and variance in square millimetres; the remaining evidence maps are
dimensionless.

\begin{figure}[t]
\centering
\includegraphics[width=\columnwidth]{figures/distrinvs_dss_evidence/distrinvs_dss_evidence.pdf}
\caption{How continuous DSS support becomes a valid warp, illustrated on a
zero-shot EndoVis dataset~8 result. (a) The completed-geometry DSS warp;
unsupported pixels are black. (b) Unnormalized support, with colors capped at
$q_{\rm vis}$ only for display; the cyan contour marks the validity threshold.
(c) The exact render-validity mask is computed from uncapped support.}
\label{fig:dss-evidence}
\end{figure}

\noindent\textbf{Depth hypotheses.}
For each source pixel $q$, depth hypothesis $z_{qk}$ corresponds to a fixed
standardized offset $\xi_k\in\Xi$ and quadrature weight $a_k$. The set and
weights are specified in Sec.~\ref{sec:implementation}. The calibrated pinhole
operator projects each hypothesis to a target coordinate
$\mathbf x_{qk}\in\R^2$ and target depth $z^t_{qk}\in\R_+$. It comprises source
back-projection, rigid transformation, and target projection
~\cite{wiles2020synsin,niklaus2020softmax}.

\noindent\textbf{Uncertainty-aware visibility.}
Depth uncertainty is converted into a hypothesis-dependent, axis-aligned
target footprint. The projection derivative $J^\alpha_{qk}$ measures the
sensitivity of image coordinate $\alpha\in\{u,v\}$ to source depth. It has
units of pixels per millimetre. The pixel-space footprint is
\begin{equation}
 r^\alpha_{qk}
 =b+|J^\alpha_{qk}|\sigma_q+b_n(1-|n_q^z|).
\label{eq:dss-footprint}
\end{equation}
The radii $r^\alpha_{qk}$, $b$, and $b_n$ are measured in pixels, while $n_q^z$
is the optical-axis component of the orientation proxy. For target pixel $p$,
$\rho_{qkp}$ denotes its elliptical distance from the projected hypothesis,
normalized by both radii. The corresponding contribution is
\begin{equation}
 w_{qk\rightarrow p}=c_q a_k\exp\!\left[
 -\frac{\rho_{qkp}^2}{2}
 -\frac{[z^t_{qk}-z_{\min}(p)]_+}{\tau_z}\right].
\label{eq:dss-weight}
\end{equation}
Source confidence and quadrature mass scale this contribution. Spatial
proximity and soft nearest-surface visibility form the exponential
penalties. Here, $z_{\min}(p)$ is the detached nearest arrival,
$[\,\cdot\,]_+$ is the positive part, and $\tau_z$ is the visibility
temperature in millimetres. Nearest-depth selection and integer assignment are
discrete; all remaining operations are differentiable.

\noindent\textbf{Reliability statistics.}
Let $\Omega_p$ collect the retained source--hypothesis contributions to target
pixel $p$. Their total weight defines raw support $S_p$. At render-valid pixels,
DSS normalizes these weights over $\Omega_p$. Their weighted moments yield the
transported tensor, expected depth, and conditional depth variance. Coverage
is a bounded saturation of raw support. The normalization constant and exact
coverage transform are reported in Sec.~\ref{sec:implementation}.

DSS couples normalized weight entropy to depth dispersion. Let $H_N(p)$ denote
the normalized Shannon entropy of the retained weights at pixel $p$. Collision
ambiguity is
\begin{equation}
 \mathcal H_p
 =H_N(p)\left[1-\exp(-V_p/\tau_z^2)\right].
\label{eq:collision-entropy}
\end{equation}
This coupling discounts entropy caused only by benign subpixel overlap. We
derive target confidence $\kappa_p$ by penalizing coverage with depth variance
and collision ambiguity. A pixel is render-valid when its raw support reaches
$S_{\min}$, as shown in Fig.~\ref{fig:dss-evidence}. Invalid pixels retain
support and coverage but zero the transported outputs and ambiguity maps. The
continuous evidence therefore distinguishes weak support from conflicting depth
order rather than reducing reliability to a binary mask.

\subsection{Reliability-Guided Expert Fusion}

\noindent\textbf{Physics-guided routing.}
Both experts receive transported evidence and a condition tensor
$\mathbf h\in\R^{6\times H\times W}$. Its channels are coverage, confidence,
normalized variance, collision entropy, normalized depth, and DSS validity.
An analytic prior $g^0\in[0,1]^{1\times H\times W}$ favours synthesis as
geometric reliability decreases and equals one outside DSS support:
\begin{equation}
\begin{gathered}
 \phi_p=\alpha_\kappa(1-\kappa_p)+\alpha_V\bar V_p
 +\alpha_H\mathcal H_p,\\
 g_p^0=\begin{cases}
 \operatorname{clip}_{[0,1]}(\phi_p),&M_p^{\rm dss}=1,\\
 1,&M_p^{\rm dss}=0.
 \end{cases}
\end{gathered}
\label{eq:physics-prior}
\end{equation}
The provisional synthesis probability is
\begin{equation}
 \widetilde g_p=\operatorname{sigmoid}\!\left[
 \ell_p^0+\eta_g\tanh R_\theta(\mathbf h)_p\right].
\label{eq:router}
\end{equation}
Here, $\ell^0=\operatorname{sg}[\operatorname{logit}(
\operatorname{clip}_{[\epsilon_g,1-\epsilon_g]}g^0)]$, and the residual router
$R_\theta$ outputs a scalar target-grid map. The factor $\eta_g$ bounds its
correction. A binary routing mask $M^{\rm rel}$ accepts DSS evidence only
where the render is valid, $S_p\geq S_{\rm rel}$, and
$V_p\leq V_{\rm rel}$. The effective synthesis gate is
\begin{equation}
 g=(1-M^{\rm rel})\odot\widetilde g.
 \label{eq:effective-gate}
\end{equation}
Routing therefore retains a geometric initialization, suppresses synthesis on
reliable DSS support, and adapts continuously elsewhere. The prior coefficients
and thresholds are reported in Sec.~\ref{sec:implementation}. All three maps
have shape $1\times H\times W$.

\noindent\textbf{Complementary experts.}
The shared expert input has $(9+d_f)\times H\times W$ channels. The transport
expert predicts a bounded RGB correction to the DSS warp, yielding
$\I_t^{\rm trans}\in[0,1]^{3\times H\times W}$. It has no depth head, so its
depth remains the DSS expectation. The synthesis expert is an
uncertainty-conditioned Fourier U-Net~\cite{ronneberger2015unet} that combines
local convolution with real-FFT processing. Its outputs
$(\I_t^{\rm syn},\D_t^{\rm syn},\U_t^{\rm syn})$ have channel dimensions
$3$, $1$, and $1$ over the target grid. This local--global design follows
Fourier inpainting networks
~\cite{suvorov2022lama,chu2023uffc}, but all weights are trained within
DistriNVS. Risk is directly trained as an error-scale proxy on unreliable
regions, not as a calibrated failure probability over the full image.

\noindent\textbf{Reliability-constrained output.}
The transport branch retains source evidence while allowing bounded appearance
correction; the synthesis branch supplies content where geometric evidence is
weak. Their final predictions are
\begin{align}
 \widehat\I_t&=(1-g)\odot\I_t^{\rm trans}+g\odot\I_t^{\rm syn},\\
 \widehat\D_t&=(1-g)\odot\bar D_t+g\odot\D_t^{\rm syn},
 \qquad \U_t=\U_t^{\rm syn}.
 \label{eq:expert-fusion}
\end{align}
The gate is broadcast across RGB channels. On $M^{\rm rel}$, $g=0$ and the
output follows the bounded transport expert rather than an exact copy of the
warp; elsewhere the two experts form a soft mixture. For regional supervision,
we open the unreliable complement and define
$M^{\rm reg}=M^{\rm dss}\land\neg\operatorname{Open}_{k}(1-M^{\rm rel})$,
where $M^{\rm dss}$ is render validity. This cleanup removes isolated
unreliable islands from the regional partition; it affects only losses and
diagnostics, never the predicted RGB-D or the gate in
Eq.~\eqref{eq:effective-gate}.

\subsection{Learning Objectives}

The training target $\I_t\in[0,1]^{3\times H\times W}$ is accompanied by a
tissue-valid mask $M_{\rm tis}\in\{0,1\}^{1\times H\times W}$, which excludes
annotated instruments when available. The unreliable, reliable, and seam
regions are $M_u=(1-M^{\rm reg})\land M_{\rm tis}$,
$M_v=M^{\rm reg}\land M_{\rm tis}$, and
$M_b=\operatorname{Band}_{w}(1-M^{\rm reg})\land M_{\rm tis}$, where
$\operatorname{Band}_{w}$ is a width-$w$ dilation--erosion band. Their
Charbonnier terms compare predicted and target RGB;
$\mathcal L_\nabla$ additionally matches finite-difference image gradients on
$M_u\cup M_b$. Training uses this cleaned regional partition, whereas regional
image-quality and risk metrics use the raw DSS hole mask.

A masked Laplace negative log-likelihood $\mathcal L_{d\text{-nll}}$ supervises
the source depth mean and scale on the pre-dropout mask $\M_s^{\rm gt}$.
Available target depth $\D_t$ shares the dimensions of $\widehat\D_t$ and
supplies a masked relative-depth term $\mathcal L_{t\text{-depth}}$. The
depth-guided cycle splats the predicted RGB-D back through the inverse calibrated
transform. Detached risk controls its depth scale, whereas complementary risk
controls its confidence. The cycle term
$\mathcal L_{\rm cyc}$ compares backward-rendered RGB with the source image
where both are valid. We do not apply a depth-cycle loss.

The detached router target is
$y_p=\mathbf{1}\{S_p<S_{\rm rel}\,\lor\,V_p>V_{\rm rel}\,\lor\,
\mathcal H_p>\mathcal H_{\rm rel}\}$, and $\mathcal L_{\rm router}$ applies
binary cross-entropy to the provisional logit in Eq.~\eqref{eq:router} over
the full image. Render-invalid pixels
necessarily satisfy the support test because $S_{\rm rel}>S_{\min}$. For a
mask $M$, $\langle\cdot\rangle_M$ denotes the mean over its selected entries.
The scalar $e_p$ is the channel-mean absolute RGB error at pixel $p$. The
positive scale $s_p$ is an affine mapping of predicted risk between fixed
bounds. The uncertainty objective is
\begin{equation}
 \mathcal L_u
 =\left\langle e_p/s_p+\log s_p\right\rangle_{M_u}.
\label{eq:risk-loss}
\end{equation}
Exact risk bounds are given in Sec.~\ref{sec:implementation}. The target-depth
term is omitted when target depth is unavailable. The complete objective is
\begin{equation}
\begin{aligned}
 \mathcal L={}&\lambda_h\mathcal L_h+\lambda_v\mathcal L_v
 +\lambda_b\mathcal L_b+\lambda_\nabla\mathcal L_\nabla\\
 &+\lambda_d(\mathcal L_{d\text{-nll}}+\mathcal L_{t\text{-depth}})
 +\lambda_c\mathcal L_{\rm cyc}\\
 &+\lambda_r\mathcal L_{\rm router}+\lambda_u\mathcal L_u.
\end{aligned}
\label{eq:total-loss}
\end{equation}

At inference, the predictor consumes only source RGB-D, camera calibration,
and relative pose, and ends at the fused prediction in
Eq.~\eqref{eq:expert-fusion}. Target RGB-D and instrument masks are used only
for supervision or evaluation; the backward-cycle branch and deterministic
image repair are disabled. The method assumes known
zero-skew pinhole intrinsics, a rigid inter-endoscope transform, and an
axis-aligned approximation to projected uncertainty; these assumptions delimit
its present scope.

\section{Experiments}
\label{sec:experiments}

\subsection{Datasets and Zero-Shot Protocol}

Training uses the frozen list of 13 iMED Task~2 scene directories supplied in
the project configuration~\cite{imed2026}. Splits are made at scene level, not
at frame level. The training loader uses a synchronized source RGB-depth and
target RGB pair at $640\times512$ resolution. Fifteen percent of valid source
depth samples are removed online, with additional $9\times9$ missing islands,
so that the reliability encoder observes both isolated and contiguous depth
failures. EndoVis/SCARED datasets 8 and 9 are reserved for final zero-shot
evaluation; neither their images nor metrics are used for model selection.

We determine inclusion using a renderer-only pass from the \emph{raw} valid
source depth. A frame is evaluated only if
$|M_{\rm raw}^{\rm dss}|/(H_IW_I)>0.5$; the learned depth completion, target image,
and target tool mask cannot influence this decision. This reproduces the
project's requested strict overlap criterion and prevents a learned module from
selecting favourable frames. The current tables report frame-wise means and
sample standard deviations, both rounded to two decimal places. Paired
sequence-level confidence intervals and permutation tests will be added once
all baseline predictions are frozen.

\todo{Verify the final dataset provenance, license statement, total training
pair count, held-out sequence count, and the immutable list of selected test
frames before submission.}

\subsection{Baselines, Metrics, and Implementation}
\label{sec:implementation}

Geometry baselines will include hard point projection, nearest-neighbour
small-hole filling, deterministic soft splatting, and softmax splatting. Image
completion baselines will include Telea, LaMa, UFFC, MAT, Stable Diffusion 1.5
Inpainting, and the existing LCM+endoscopy-LoRA system. Directly relevant
endoscopic baselines (DAEVI and SSIFNet) will be evaluated when code and an
equivalent cross-view protocol are available; otherwise their non-comparable
reported values will not be inserted into the main table. Every completion
method will use the same raw mask, spatial resolution, selected frames, and
direct-output evaluation without method-specific image repair.

The existing evaluator reports full-, raw-hole-, raw-visible-, and morphological
seam-band PSNR/SSIM, full-image and hole-region LPIPS, known-region drift, risk
expected calibration error (ECE), area under the risk--coverage curve (AURC),
and a projection-defect taxonomy. The final study will additionally report DISTS, depth AbsRel/RMSE,
disparity endpoint error, normal error, latency, throughput, peak memory, and
parameter count. Results will be stratified by raw overlap, hole ratio, depth
validity, instrument presence, and sequence.

\noindent\textbf{Renderer and reliability.}
We use $\Xi=(-1,0,1)$ and $(a_k)=(0.2741,0.4518,0.2741)$; the display
percentile in Fig.~\ref{fig:dss-evidence} is $q_{\rm vis}=99$ and affects
visualization only. Before projection, every sampled depth
hypothesis $z_{qk}$ is clipped to $[1,500]$ mm. Hypotheses affect the radius-one
integer neighbourhood, with $\tau_z=2$ mm and $S_{\min}=0.03$. Each footprint
$r^\alpha_{qk}$ is clipped to $[0.35,2.5]$ px, with $b=b_n=0.75$ px. Spatial
and final weights no greater than $10^{-5}$ and $10^{-8}$ are pruned, and
$\epsilon_{\rm dss}=10^{-8}$ stabilizes normalization.
Coverage uses $C=1-\exp(-S)$, and target confidence uses
$\kappa=C\exp[-V/(4\tau_z^2)-\mathcal H]$. Weight entropy is divided by the
logarithm of at least two retained contributions before it is coupled to depth
variance.

The reliability encoder fuses full-resolution residual and quarter-resolution
context branches to output $d_f=16$ feature channels. Its fallback median is
$s_D=100$ mm, and valid depth is normalized as
$\widetilde D_s=\operatorname{clip}(D_s/s_D,0,5)$ before missing values are set
to zero. Measured and completed confidence occupy $[0.5,1]$ and $[0,0.35]$,
respectively. The measured-depth correction ratio is bounded by $0.1$, whereas
the completion branch has a $0.05$ positive offset. Depth and uncertainty are
clipped to $[1,500]$ mm and $[0.25,20]$ mm.

\noindent\textbf{Router and experts.}
The condition maps use normalized variance
$\bar V=V/(V+\tau_z^2)$ and normalized depth
$\widetilde D=\operatorname{clip}(\bar D_t/s_D,0,5)/5$. The analytic prior is
clipped to $[0,1]$ and then to $[\epsilon_g,1-\epsilon_g]$ before its detached
logit $\ell^0$ is evaluated. For Eq.~\eqref{eq:router}, the prior weights are
$(\alpha_\kappa,\alpha_V,\alpha_H)=(0.65,0.20,0.15)$,
$\epsilon_g=10^{-4}$, and $\eta_g=0.25$; the residual router is
zero-initialized. The transport RGB residual is bounded by $0.15$. The
four-resolution synthesis network uses $3\times3$ local branches and five
output channels; its metric depth is
$[\operatorname{softplus}(\cdot)+0.05]s_D$. The shared base width is 32.

\noindent\textbf{Routing masks and losses.}
The routing tests use $S_{\rm rel}=0.35$, $V_{\rm rel}=9$ mm$^2$, and
$\mathcal H_{\rm rel}=1.0$; regional-mask cleanup uses a $7\times7$ opening kernel. Since
collision entropy lies in $[0,1]$, its configured hard threshold is neutral,
although entropy still affects target confidence and soft routing. A separate
$7\times7$ dilation--erosion kernel extracts the training seam. The
Charbonnier constant is $10^{-3}$, the source-depth likelihood uses a
$10^{-3}$ scale floor, and target relative-depth error uses a 1-mm denominator
floor. Equation~\eqref{eq:risk-loss} uses
$(s_{\min},s_{\max})=(0.01,0.50)$ in normalized-RGB units. For the backward
cycle, the depth scale is
$\operatorname{clip}_{[0.25,20]}[0.25+20\operatorname{sg}(\U_t)]$ mm and its
confidence is $\operatorname{clip}_{[0.05,1]}(1-\U_t)$.

\noindent\textbf{Optimization.}
Training is scheduled for 10,000 AdamW steps with learning rate
$2\times10^{-4}$, weight decay $10^{-4}$, and three-GPU distributed training.
The per-GPU batch size is 1 with four-step gradient accumulation, giving an
effective batch size of 12. We use BF16 mixed precision, maximum gradient norm
1, and seed 6666. The implemented loss weights
$(\lambda_h,\lambda_v,\lambda_b,\lambda_\nabla,\lambda_d,\lambda_c,
\lambda_r,\lambda_u)$ are $(3,1,2,0.2,0.2,0.2,0.1,0.1)$.

\todo{Record the final checkpoint, model-selection rule, GPU model, training
time, inference warm-up, precision, exact parameter/FLOP counts, and all
baseline versions and checkpoints.}

\subsection{Main Comparison}

We evaluated every available prediction on the seven-scene iMED validation
split. Complete-coverage methods produced predictions for all 1,392 frames,
and every included prediction was re-evaluated against the same target image
and raw-warp hole mask. The mask covered 25.84\% of an image on average.
Five Gaussian pipelines covered all seven scenes. The four published
reconstruction baselines used dataset-provided ground-truth depth, whereas the
additional StereoSurGS pipeline combined SurgicalGS with
FoundationStereo-estimated depth~\cite{chen2025surgicalgs,wen2025foundationstereo}.
They are reference reconstruction baselines rather than matched feed-forward
completion models.

PR-ENDO, EndoGS, and EndoGaussian met the original run-validity criterion on
1/7, 5/7, and 6/7 scenes (199, 994, and 1,193 frames), respectively. We report
these method-dependent subsets as diagnostics, but exclude them from ranking
and best-value selection. The EndoGaussian diagnostic retains all frames from
its six accepted scenes, including 155 frames flagged as degenerate, rather
than selectively removing failed frames. D4Recon yielded only degenerate
renders on all 1,392 selected frames, while the remaining zero-success runs
provided no evaluable render set; assigning either group an image-quality
score would therefore be artificial. Table~\ref{tab:validation-comparison}
reports the frame-wise mean and sample standard deviation under this protocol.
The DistriNVS row evaluates the primary soft-fusion output directly, without
output hard copying or small-component repair.

\begin{table*}[t]
\centering
\caption{Comparison on the seven-scene iMED validation split. Rows with
$N=1{,}392$ cover all seven scenes; partial rows summarize only their accepted
scenes and are diagnostic. Every included prediction uses the same target and
raw-warp hole mask. Values are frame-wise mean~$\pm$~sample SD. Bold denotes
the best mean among complete-coverage rows only. FS denotes
FoundationStereo-estimated depth; the other Gaussian runs use dataset
ground-truth depth. DistriNVS uses direct soft fusion without output hard
copying or small-component repair.}
\label{tab:validation-comparison}
\small
\setlength{\tabcolsep}{2.0pt}
\begin{tabular}{lccccccc}
\toprule
Method & $N$ & PSNR $\uparrow$ & SSIM $\uparrow$ & LPIPS $\downarrow$ &
H-PSNR $\uparrow$ & H-SSIM $\uparrow$ & H-LPIPS $\downarrow$ \\
\midrule
\multicolumn{8}{l}{\textit{Gaussian reconstruction---complete coverage}} \\
Deform3DGS~\cite{yang2024deform3dgs} & 1,392 & $15.55\pm1.78$ & $0.55\pm0.04$ &
$0.63\pm0.06$ & $14.11\pm1.37$ & $0.47\pm0.10$ & $0.69\pm0.06$ \\
Endo-4DGS~\cite{huang2024endo4dgs} & 1,392 & $16.44\pm1.71$ & $0.58\pm0.04$ &
$0.61\pm0.06$ & $15.61\pm1.47$ & $0.57\pm0.13$ & $0.66\pm0.06$ \\
Free-SurGS~\cite{guo2024freesurgs} & 1,392 & $16.83\pm2.11$ & $\mathbf{0.67\pm0.06}$ &
$0.63\pm0.04$ & $17.05\pm2.38$ & $0.72\pm0.04$ & $0.61\pm0.04$ \\
SurgicalGS~\cite{chen2025surgicalgs} & 1,392 & $17.41\pm1.99$ & $0.60\pm0.06$ &
$0.60\pm0.05$ & $16.98\pm2.55$ & $0.61\pm0.14$ & $0.65\pm0.04$ \\
StereoSurGS (FS) & 1,392 & $16.37\pm1.46$ & $0.52\pm0.08$ &
$0.59\pm0.05$ & $13.59\pm1.27$ & $0.39\pm0.13$ & $0.70\pm0.06$ \\
\addlinespace
\multicolumn{8}{l}{\textit{Gaussian reconstruction---partial coverage (diagnostic)}} \\
PR-ENDO~\cite{kaleta2025prendo} & 199 & $15.51\pm0.76$ & $0.50\pm0.01$ &
$0.62\pm0.01$ & $12.98\pm0.46$ & $0.20\pm0.05$ & $0.71\pm0.02$ \\
EndoGS~\cite{zhu2024endogs} & 994 & $15.86\pm2.19$ & $0.61\pm0.06$ &
$0.71\pm0.11$ & $15.50\pm1.97$ & $0.65\pm0.08$ & $0.70\pm0.07$ \\
EndoGaussian~\cite{liu2025endogaussian} & 1,193 & $15.21\pm3.56$ &
$0.57\pm0.26$ & $0.81\pm0.10$ & $16.00\pm3.15$ & $0.62\pm0.28$ &
$0.72\pm0.08$ \\
\addlinespace
\multicolumn{8}{l}{\textit{Image completion}} \\
SD1.5 Inpainting & 1,392 & $16.84\pm2.73$ & $0.57\pm0.05$ & $0.52\pm0.06$ &
$14.56\pm3.14$ & $0.50\pm0.10$ & $0.64\pm0.07$ \\
LaMa & 1,392 & $18.99\pm1.58$ & $0.64\pm0.07$ & $0.47\pm0.03$ &
$20.41\pm2.06$ & $0.74\pm0.05$ & $0.48\pm0.04$ \\
LCM + endoscopy LoRA & 1,392 & $19.02\pm1.64$ & $0.63\pm0.09$ & $0.49\pm0.03$ &
$19.54\pm2.28$ & $0.74\pm0.09$ & $0.49\pm0.03$ \\
\midrule
DistriNVS (soft fusion; primary) & 1,392 & $\mathbf{22.95\pm2.33}$ & $0.59\pm0.11$ &
$\mathbf{0.32\pm0.05}$ & $\mathbf{27.95\pm2.91}$ & $\mathbf{0.80\pm0.07}$ &
$\mathbf{0.28\pm0.06}$ \\
\bottomrule
\end{tabular}
\end{table*}

The primary soft-fusion DistriNVS output achieved the best PSNR and LPIPS
among complete-coverage methods. Relative to the strongest image-completion
baselines, it improved full-image PSNR by 3.93~dB over LCM and reduced LPIPS
by 0.15 relative to LaMa. Inside the raw DSS holes, it improved PSNR by
7.54~dB, SSIM by 0.06, and LPIPS by 0.20 over the strongest corresponding
completion baselines. Full-image SSIM was lower than LaMa by 0.05 and
Free-SurGS by 0.08, indicating that the gains in pixel fidelity and perceptual
distance do not translate into the best global structural score. The
partial-coverage diagnostics characterize different subsets and do not
support pairwise superiority claims.

We separately evaluated DistriNVS on held-out EndoVis Datasets~8 and 9 using
the strict $\mathrm{raw\ overlap}>0.5$ inclusion rule. Among 124 candidate
frames, 78 met the rule. Table~\ref{tab:zeroshot-breakdown} reports the two
datasets separately and jointly.

\begin{table*}[t]
\centering
\caption{Legacy DistriNVS zero-shot diagnostic on held-out EndoVis Datasets~8 and 9.
Values are frame-wise mean~$\pm$~SD over frames satisfying
$\mathrm{raw\ overlap}>0.5$. This archived run used output hard copying and
small-component repair and will be replaced by direct-output results.}
\label{tab:zeroshot-breakdown}
\setlength{\tabcolsep}{3.2pt}
\begin{tabular}{lrrrrrrr}
\toprule
Dataset & Frames & PSNR $\uparrow$ & SSIM $\uparrow$ & LPIPS $\downarrow$ &
H-PSNR $\uparrow$ & H-SSIM $\uparrow$ & H-LPIPS $\downarrow$ \\
\midrule
EndoVis 8 & 48 & $\mathbf{24.85\pm2.00}$ & $\mathbf{0.72\pm0.12}$ &
$\mathbf{0.21\pm0.04}$ & $\mathbf{23.71\pm2.74}$ &
$\mathbf{0.70\pm0.09}$ & $\mathbf{0.23\pm0.06}$ \\
EndoVis 9 & 30 & $22.10\pm2.29$ & $0.63\pm0.09$ & $0.27\pm0.06$ &
$18.80\pm3.32$ & $0.45\pm0.08$ & $0.35\pm0.08$ \\
\midrule
Combined & 78 & $23.79\pm2.49$ & $0.68\pm0.12$ & $0.23\pm0.06$ &
$21.82\pm3.81$ & $0.61\pm0.15$ & $0.28\pm0.09$ \\
\bottomrule
\end{tabular}
\end{table*}

Across the 78 accepted frames, the archived variant obtained
$23.79\pm2.49$~dB PSNR, $0.68\pm0.12$ SSIM,
$0.23\pm0.06$ LPIPS, $21.82\pm3.81$~dB hole PSNR,
$0.61\pm0.15$ hole SSIM, and $0.28\pm0.09$ hole LPIPS. Baseline inference on
Datasets~8 and 9 is not yet complete; consequently,
Table~\ref{tab:zeroshot-breakdown} does not support a cross-method zero-shot
ranking. Paired confidence intervals and significance tests will be added after
the final prediction set is frozen.

\subsection{Ablation Study}

Table~\ref{tab:ablation} reports the nine completed configurations on the same
1,392 frames from the seven-scene iMED validation split. All models were
trained for 10,000 steps with seed 6666 and evaluated without small-component
repair. Values in the table are frame-wise mean~$\pm$~sample SD, and paired
95\% confidence intervals below were obtained by resampling the seven scenes.
The primary method is the soft fusion (SF) configuration in
Eq.~\eqref{eq:expert-fusion}. The seven completed component variants retain
legacy hard composition (HC) and are therefore compared only with the HC
reference, rather than treated as removals from the primary model. For compact
ablation labels, we use depth completion (DC), distributional depth (DD),
anisotropic footprints (AF), physics-guided routing (PGR), RGB-D cycle
consistency (RCC), uncertainty loss (UL), and trust-mask cleanup (TMC).

\begin{table*}[t]
\centering
\caption{Single-seed ablation on the seven-scene iMED validation split
($N=1{,}392$ frames). H-PSNR and H-SSIM use each configuration's raw DSS hole
mask; drift is the mean absolute RGB deviation from the DSS warp on the
cleaned reliable region. Values are frame-wise mean~$\pm$~sample SD from direct
outputs without small-component repair. Bold marks the best unrounded
image-quality or risk mean; hard-copy variants have zero drift by construction.}
\label{tab:ablation}
\small
\setlength{\tabcolsep}{1.5pt}
\begin{tabular}{lrrrrrrr}
\toprule
Variant & PSNR $\uparrow$ & SSIM $\uparrow$ & H-PSNR $\uparrow$ &
H-SSIM $\uparrow$ & Seam PSNR $\uparrow$ & Drift $\downarrow$ & AURC $\downarrow$ \\
\midrule
\multicolumn{8}{l}{\textit{Primary output rule and legacy reference}} \\
DistriNVS (SF; primary) & $\mathbf{22.95\pm2.33}$ & $\mathbf{0.59\pm0.11}$ &
$\mathbf{27.95\pm2.91}$ & $\mathbf{0.80\pm0.07}$ & $\mathbf{24.56\pm2.33}$ & $0.05\pm0.01$ &
$\mathbf{0.02\pm0.01}$ \\
DistriNVS (HC; reference) & $20.75\pm1.90$ & $0.49\pm0.13$ & $27.56\pm2.78$ &
$0.78\pm0.07$ & $22.06\pm1.70$ & $0.00\pm0.00$ & $0.03\pm0.01$ \\
\addlinespace
\multicolumn{8}{l}{\textit{Component variants under legacy hard-copy output}} \\
w/o DC & $20.68\pm1.89$ & $0.49\pm0.13$ & $27.76\pm2.48$ &
$0.79\pm0.06$ & $22.11\pm1.66$ & $0.00\pm0.00$ & $0.03\pm0.01$ \\
w/o DD ($K=1$) & $20.67\pm1.91$ & $0.48\pm0.13$ & $27.19\pm2.74$ &
$0.77\pm0.07$ & $21.51\pm1.66$ & $0.00\pm0.00$ & $0.03\pm0.01$ \\
w/o AF & $20.65\pm1.92$ & $0.48\pm0.13$ & $27.73\pm2.71$ &
$0.78\pm0.06$ & $22.11\pm1.73$ & $0.00\pm0.00$ & $0.02\pm0.01$ \\
w/o PGR & $20.75\pm1.89$ & $0.49\pm0.13$ & $27.41\pm2.72$ &
$0.77\pm0.07$ & $21.82\pm1.68$ & $0.00\pm0.00$ & $0.03\pm0.01$ \\
w/o RCC & $20.72\pm1.91$ & $0.49\pm0.13$ & $27.41\pm2.73$ &
$0.78\pm0.07$ & $21.98\pm1.70$ & $0.00\pm0.00$ & $0.03\pm0.01$ \\
w/o UL & $20.70\pm1.89$ & $0.48\pm0.13$ & $27.02\pm2.60$ &
$0.77\pm0.07$ & $21.96\pm1.71$ & $0.00\pm0.00$ & $0.04\pm0.01$ \\
w/o TMC & $20.77\pm1.91$ & $0.49\pm0.13$ & $27.56\pm2.73$ &
$0.78\pm0.07$ & $22.39\pm1.69$ & $0.00\pm0.00$ & $0.03\pm0.01$ \\
\bottomrule
\end{tabular}
\end{table*}

The primary soft-fusion configuration achieved 22.949~dB full-image PSNR,
0.589 SSIM, 27.951~dB hole PSNR, 0.803 hole SSIM, and 24.563~dB seam
PSNR. Relative to the legacy hard-copy reference, the paired differences were
$+2.202$~dB for full-image PSNR (95\% CI, $+1.649$ to $+2.799$) and
$+2.504$~dB for seam PSNR (95\% CI, $+1.512$ to $+3.612$). AURC decreased
from 0.0269 to 0.0206, while reliable-region drift increased from zero to
0.047 (paired 95\% CI, $+0.042$ to $+0.052$) in normalized RGB units. This
quality--fidelity trade-off supports soft fusion as the primary output rule.

Within the legacy hard-copy family, replacing distributional depth with a
single deterministic sample produced the clearest consistent degradation:
full-image, hole, and seam PSNR decreased by 0.075~dB (95\% CI, $-0.107$ to
$-0.048$), 0.367~dB (95\% CI, $-0.637$ to $-0.177$), and 0.546~dB (95\%
CI, $-0.801$ to $-0.311$), respectively. Removing the physics router left
full-image PSNR nearly unchanged ($-0.003$~dB; 95\% CI, $-0.012$ to
$+0.007$) but reduced seam PSNR by 0.241~dB (95\% CI, $-0.319$ to
$-0.158$). Depth completion and anisotropic footprints showed mixed regional
effects: their removal reduced full-image PSNR by 0.069 and 0.095~dB,
respectively, but changed hole PSNR by $+0.203$~dB (95\% CI, $-0.109$ to
$+0.525$) and $+0.170$~dB (95\% CI, $+0.056$ to $+0.313$). The deterministic
renderer also reduced mean raw overlap from 0.765 to 0.758, so its hole and
seam differences include a modest change in the evaluated masks.

Disabling trust-mask cleanup left hole PSNR and hole SSIM effectively
unchanged, while increasing seam PSNR by 0.329~dB (95\% CI, $+0.206$ to
$+0.458$) and reducing AURC from 0.0269 to 0.0254. The full-image PSNR change
was small and inconclusive ($+0.019$~dB; 95\% CI, $-0.005$ to $+0.044$),
whereas SSIM decreased by 0.0013 (95\% CI, $-0.0018$ to $-0.0008$). Cleanup
increased the trusted hard-copy ratio from 0.755 to 0.762, but did not improve
aggregate seam fidelity under this single-seed evaluation.

Removing the RGB-D cycle caused a small but consistent reduction in
full-image PSNR ($-0.028$~dB; 95\% CI, $-0.045$ to $-0.008$) and seam PSNR
($-0.083$~dB; 95\% CI, $-0.122$ to $-0.039$). Its hole-PSNR change was less
conclusive ($-0.148$~dB; 95\% CI, $-0.364$ to $+0.061$), indicating that the
cycle objective primarily provided a modest regularization benefit. Removing
the uncertainty loss reduced hole PSNR by 0.541~dB (95\% CI, $-0.857$ to
$-0.255$) and seam PSNR by 0.102~dB (95\% CI, $-0.171$ to $-0.038$). More
importantly, risk AURC increased from 0.0269 to 0.0365, while risk ECE
increased from 0.0236 to 0.9421. Thus, within the legacy hard-copy family,
the uncertainty objective was essential for calibrated risk estimates and
also improved reconstruction quality in the raw DSS holes.

\subsection{Qualitative, Failure, and Efficiency Analysis}

\begin{figure*}[t]
\centering
\fbox{\parbox[c][35mm][c]{0.96\textwidth}{\centering
\todo{Replace this placeholder with synchronized columns: source RGB, raw hard
warp, raw validity, LaMa, LCM, DistriNVS, target RGB, support, variance,
collision entropy, and risk. Include at least one success and one failure for
each overlap bin.}}}
\caption{Planned qualitative comparison and failure visualization. Crops and
color scales must be identical across methods; generated and source-backed
regions will be explicitly marked.}
\label{fig:qualitative}
\end{figure*}

Efficiency is central to the ICRA framing. We will measure end-to-end latency
including depth reliability estimation, DSS, routing, and both experts;
DSS-only timing will also be reported to isolate the renderer cost. Failure
analysis will cover specular depth loss, calibration perturbation, thin
instruments, large disocclusion, weak texture, and dynamic tissue deformation.

\todo{Insert the speed/memory/parameter table, calibration-perturbation curves,
and blinded expert assessment if available. Do not claim real-time performance
until measured on the stated hardware and resolution.}

\section{Discussion}
\label{sec:discussion}

The ablation separates two complementary design effects. Within the legacy
hard-copy family, distributional depth consistently improved full-image, hole,
and seam fidelity, whereas the physics router primarily improved seam quality.
Depth completion and the anisotropic footprint produced smaller, region-
dependent changes. Trust-mask cleanup slightly increased global SSIM and the
trusted hard-copy ratio, whereas disabling it improved seam PSNR and AURC.
The output rule yielded the largest difference: the
separately trained primary soft-fusion configuration improved full-image and
seam PSNR by 2.20 and 2.50~dB and reduced AURC from 0.0269 to 0.0206 relative
to the legacy hard-copy reference, at the cost of 0.047 reliable-region drift.
This pattern is consistent with hard copying carrying residual transport error
to the output, whereas soft fusion permits bounded appearance correction.
Because the two output configurations differ during both training and
inference, the comparison does not isolate these two effects. The suite also
contains one seed, and its component variants were trained under the legacy
output rule; their interactions with the primary soft-fusion configuration are
therefore not yet established. Moreover, the zero-shot table still uses the
old hard-copy/repair protocol; cross-dataset conclusions must be refreshed
from direct no-post outputs before submission.

The method does not remove the fundamental ambiguity of unseen tissue. A low
risk score indicates model confidence under the training distribution; it does
not certify anatomical truth. Moreover, the current formulation assumes known
intrinsics and a rigid inter-camera transform for each pair, processes frames
independently, and uses an axis-aligned rather than full-covariance footprint.
The archived zero-shot run also varied substantially between the two held-out
datasets. Because that run used the legacy output protocol and baseline
coverage remains incomplete, it cannot establish uniform cross-dataset
generalization or a fair ranking on Datasets~8 and 9.
Large calibration errors, non-rigid motion between nominally synchronized
frames, or out-of-distribution anatomy may therefore invalidate the transport
statistics. Generated output should remain visibly marked and should not be
used as a substitute for a live physical endoscopic view.

\section{Conclusion}
\label{sec:conclusion}

This paper introduces a reliability-aware formulation of cross-endoscope
novel-view synthesis and implements it through distributional surface
splatting, reliability-constrained transport/synthesis fusion, and
joint appearance--geometry risk modeling. The design
turns projection support, depth variance, and collision entropy into operational
constraints on what the network may synthesize. No-post validation on the
seven-scene iMED split showed that the primary soft-fusion configuration
improved full-image and seam PSNR by 2.20 and 2.50~dB over the legacy hard-copy
reference. Within the legacy component study, distributional depth provided
the most consistent image-fidelity benefit, while physics routing primarily
improved transition boundaries. These results support direct reliability-
constrained fusion as the primary output formulation. The current evidence is
limited to one seed, with component variants evaluated under the legacy
hard-copy rule; matched multi-seed and held-out zero-shot experiments remain
necessary for broader comparative claims.

\bibliographystyle{IEEEtran}
\bibliography{references}

\end{document}
